\documentclass[10pt,a4paper]{amsart}
\usepackage[margin=19mm]{geometry}
\usepackage{amsmath,amssymb,mathtools}
\usepackage[scr=rsfs,cal=euler]{mathalfa}
\usepackage{xfrac}
\usepackage[unicode,hidelinks,bookmarks=false]{hyperref}
\newcommand{\ie}{i.e.\ }
\newcommand{\cf}{cf.\ }
\newcommand{\resp}{respectively\ }
\newcommand{\Subsection}{\subsection}
\providecommand{\nobreakdash}{\nobreak\hbox{-}}
\setlength{\emergencystretch}{2em}
\multlinegap0pt
\usepackage{amssymb}
\usepackage{mathtools}
\usepackage[shortlabels]{enumitem}
\usepackage{xcolor}
\usepackage{bbm}
\newtheorem{lemma}{Lemma}
\newcommand{\Cstar}{\ensuremath{C^*}}
\newcommand{\Eboundary}{E^0_{\mathrm{bdry}}}
\newcommand{\Eint}{E^0_{\mathrm{int}}}
\newcommand{\edgebdy}{E^1_{\mathrm{bdry}}}
\newcommand{\edgeint}{E^1_{\mathrm{int}}}
\newcommand{\norm}[1]{\left\|#1\right\|} % Norm
\newcommand\numberthis{\addtocounter{equation}{1}\tag{\theequation}}
\title{Graph C*-algebras are singly generated}
\author{Jakub Curda, Julian Gonzales, Victor Wu}
\date{Source: arXiv:2601.01249v1 (CC BY 4.0)}
\begin{document}
\maketitle

\noindent\emph{Source and licence.} Graph C*-algebras are singly generated. Version arXiv:2601.01249v1 (CC BY 4.0), distributed by arXiv under the Creative Commons Attribution 4.0 International licence, http://creativecommons.org/licenses/by/4.0/, as recorded in the arXiv licence metadata for this exact version and shown on the abstract page https://arxiv.org/abs/2601.01249v1. Full licence notice: \texttt{licenses/arxiv-2601.01249-CC-BY-4.0.txt}.\par\smallskip

\noindent\textit{Editorial scope.} Counted unit: the article's lemma lemma:proposition-4-part1 (source0.tex lines 510--513). The article's complete original proof of it is delivered with it (source0.tex lines 516--562, one \texttt{proof} environment of 47 lines). No mathematics is written, repaired or re-derived: the statement and the proof are the article's own lines, in the article's own notation, and no retained line is substituted or re-flowed.  Retained uncounted as the source-local context the proof needs: lines 342--360; lines 446--448; lines 470--476.  Nothing was omitted from any retained span, so the package carries no omission record.  No retained line contains a citation, so the wrapper carries no bibliography and no bibliography entry.  No retained line contains a figure environment, an image-inclusion command or drawing code, so no image is delivered and the package declares no asset dependency.  Editorial formatting only, in addition: an AMS \texttt{amsart} preamble that restates the article's own declarations and macros used by the retained lines, namely the environments \texttt{lemma} on separate counters, together with the standard packages those lines need.  The retained environments print in this document as Lemma 1 in order of appearance, whereas the article prints the same items with the numbering of its own full document.  The complete native source of the article as distributed in the arXiv e-print for this exact version is bundled unchanged as \texttt{sources/192065/source0.tex}.
\begin{lemma}\label{lemma:orthogonality}
    We have the following relations for all integers $m,n,m',n',k \ge 1$, edges $e, e' \in E^1$, and vertices $y \in Y$:
    \begin{align*}
    s_e^* s_{e'} &= 0 &&\text{ if } e \neq e'\,,\numberthis\\
    q_m &\perp s_e  &&\text{ if } e \in \edgebdy \cup \edgeint\,,\\
    q_m &\perp p_v  &&\text{ if } v \in \Eboundary \cup \Eint\,,\\    
    p_v t_{m,n} &= 0 &&\text{ if } v \in \Eboundary \cup \Eint\,,\\
    t_{m,n} t_{m',n'} &= 0\,, \\
    s_{y,k} s_{y,n}^* s_{y,n} &= s_{y,k}\,.
    \end{align*}
    Moreover, 
    \begin{equation}
        q_{m'} t_{m,n} =
    \begin{cases}
        t_{m,n} \quad &\text{if } m = m'\\
        0 \quad &\text{otherwise}
    \end{cases}\,.
    \end{equation}
\end{lemma}

\begin{equation}\label{eq:a_1}
    a_1 \coloneq \left( \sum_{y \in Y} \alpha_{y,1} s_{y,1} \right)\xi^{\frac{1}{2}} + \sum_{\substack{y \in Y \\ n \ge 1}} \alpha_{y,n+1} s_{y,n+1} s_{y,n}^*\,.
\end{equation}

\begin{lemma}\label{lemma:a_1}
    Let $a_1$ be as in \eqref{eq:a_1}. Then $a_1 = \sum_{v \in \Eboundary} p_v a_1$ as a norm convergent sum, and the inequality
    \begin{equation}\label{eq:decay-of-a_1}
    \sum_{\ell \ge m} \sum_{v \in V_\ell} \norm{p_v a_1} \le \delta_m
    \end{equation}
    holds for any $m$.
\end{lemma}

\begin{lemma}%[Proposition 4 part 1]
\label{lemma:proposition-4-part1}
    For any $M,N \ge 1$ we have $t_{M,N} \in \Cstar(g_{M,N})$ and $q_{M} \in \Cstar(g_{M+1,1} + \delta_M q_M)$.
\end{lemma}

\begin{proof}
    For $k \ge 1$, we have the following by Lemma~\ref{lemma:orthogonality}:
    \begin{align*}
    \frac{1}{\gamma_{M,N}^k}g_{M,N}^k &= \ 
    \frac{1}{\gamma_{M,N}^k}a_M^k 
    + \frac{1}{\gamma_{M,N}^k} \sum_{m \ge M} \delta_m^k q_m + \sum_{(m,n) \ge (M,N)} \frac{\gamma_{m,n}^k-\delta_m^k}{\gamma_{M,N}^k} t_{m,n}t_{m,n}^*\numberthis\\
    &+ \sum_{(m,n) \ge (M,N)} \frac{\beta_{m,n} t_{m,n} }{\gamma_{M,N}^k} \left( \sum_{j=1}^{k} \gamma_{m,n}^{k-j} \ a_M^{j-1} \right).
    \end{align*}
    By assumption we have $\delta_M < \gamma_{M,N}$, and $\norm{a_M} \le \delta_M$ by Lemma~\ref{lemma:a_1}, so the first term converges to $0$ as $k \to \infty$. The second term converges to 0 by the monotone convergence theorem, since $\delta_m \le \delta_M < \gamma_{M,N}$ for all $m \ge M$. For the third and fourth terms, the sums over $(m,n) \ge (M, N+1)$ will converge to 0 as $k \to \infty$. Indeed
    \begin{equation}
    \norm{\sum_{(m,n) \ge (M,N+1)} \frac{\gamma_{m,n}^k-\delta_m^k}{\gamma_{M,N}^k} t_{m,n}t_{m,n}^*}
    \le \max_{(m,n) \ge (M,N+1)} \left( \frac{\gamma_{m,n}^k-\delta_m^k}{\gamma_{M,N}^k}\right)
    \le \frac{\gamma_{M,N+1}^k}{\gamma_{M,N}^k}
    \end{equation}
    since the projections $t_{m,n}t_{m,n}^*$ are mutually orthogonal and the double sequence $(\gamma_{m,n})$ is decreasing with respect to the lexicographic ordering. Note that $\sum_{j=1}^\infty \gamma_{M,N+1}^{-j} \norm{a_M}^{j-1}$ converges since $\delta_M < \gamma_{M,N+1}$. Then for $(m,n) \ge (M,N+1)$ we have
    \begin{equation}
    \norm{\frac{\beta_{m,n} t_{m,n}}{\gamma_{M,N}^k} \sum_{j=1}^{k} \gamma_{m,n}^{k-j} \ a_M^{j-1}} \le \frac{\beta_{m,n} \gamma_{M,N+1}^k}{\gamma_{M,N}^k} \sum_{j=1}^\infty \gamma_{M,N+1}^{-j} \norm{a_M}^{j-1} \lesssim \beta_{m,n} \left(\frac{\gamma_{M,N+1}}{\gamma_{M,N}}\right)^k,
    \end{equation}
    and hence the sum of these terms converges to 0 as $k \to \infty$. 
    Overall, we get
    \begin{equation}\label{eq:converge-to-y}
    \frac{1}{\gamma_{M,N}^k}g_{M,N}^k \to t_{M,N}t_{M,N}^* + \beta_{M,N} t_{M,N} \sum_{j=1}^\infty \gamma_{M,N}^{-j} a_M^{j-1} \eqqcolon y
    \end{equation}
    as $k \to \infty$, so $y \in C^*(g_{M,N})$. Write
    \begin{equation}
    z \coloneq \beta_{M,N} \sum_{j=1}^\infty \gamma_{M,N}^{-j} a_M^{j-1}\,,
    \end{equation}
    so $y = t_{M,N}t_{M,N}^* + t_{M,N} z$. Then
    $
    yy^* = t_{M,N}t_{M,N}^* + t_{M,N} z z^* t_{M,N}^*\,.
    $
    Since $\norm{a_M} \le \delta_M$ and $\beta_{M,N} = \frac{1}{2} (\gamma_{M,N} - \delta_M)$, we get that
    \begin{equation}
    \norm{z} \leq \frac{\beta_{M,N}}{\gamma_{M,N} - \norm{a_M}} \leq \frac{1}{2}\,.
    \end{equation}
    Hence $yy^*$ is a small perturbation of the projection $t_{M,N}t_{M,N}^*$, so that
    \begin{equation}
    \lim_{k\to\infty} (yy^*)^{1/k} = t_{M,N}t_{M,N}^*\,.
    \end{equation}
    Thus $t_{M,N}t_{M,N}^* \in C^*(g_{M,N})$. Finally, note that
    \begin{equation}
    t_{M,N}t_{M,N}^* g_{M,N} = \beta_{M,N} t_{M,N} + \gamma_{M,N} t_{M,N}t_{M,N}^*\,
    \end{equation}
    so $t_{M,N} \in C^*(g_{M,N})$, as claimed.

    Next, consider the limit of $\frac{1}{\delta_M^k}\big(g_{M+1,1} + \delta_M q_M\big)^k$ as $k \to \infty$. The elements $g_{M+1,1}$ and $q_M$ are orthogonal by Lemma~\ref{lemma:orthogonality}, hence $\big(g_{M+1,1} + \delta_M q_M\big)^k = g_{M+1,1}^k + \delta_M^k q_M$. By \eqref{eq:converge-to-y} the limit $\lim_{k \to \infty}\frac{1}{\gamma_{M+1,1}^k}g_{M+1,1}^k$ exists, hence $\lim_{k \to \infty}\frac{1}{\delta_{M}^k}g_{M+1,1}^k = 0$ since $\gamma_{M+1,1} < \delta_M$. It follows that $\lim_{k \to \infty} \frac{1}{\delta_M^k}\big(g_{M+1,1} + \delta_M q_M\big)^k = q_M$. So indeed $q_m \in \Cstar(g_{M+1,1} + \delta_M q_M)$, as claimed.
\end{proof}

\end{document}
